\section{Innate Mars (Maurice)}
\label{app:innate_mars}

Innate Mars is Innate's Maurice mobile manipulator, integrated in \texttt{emet} for both MuJoCo simulation and ROS\,2 hardware bridges.
This appendix records embodiment details relevant to Dynagraph/DynaMem mapping and EQA---especially depth policy, which differs from Stretch and Galaxea.

\subsection{Assets and simulation}

\begin{itemize}
  \item \textbf{MJCF:} \texttt{src/emet/assets/robot/innate\_mars/innate\_mars.xml} (robot-only; room floor from \texttt{scene\_environment.xml}).
  \item \textbf{URDF reference:} \texttt{maurice.urdf} for kinematic parity with innate-os; MuJoCo adds head stereo, optical-frame corrections, and sim-only wrist camera offsets documented in \texttt{docs/robots/innate\_mars.md}.
  \item \textbf{Serve:} \texttt{emet serve mujoco --robot innate\_mars} exposes head stereo (\texttt{head\_left}, \texttt{head\_right}) and wrist \texttt{camera\_arm} over ZMQ.
  \item \textbf{Base motion:} planar slide-X / slide-Y / yaw actuators (\texttt{base\_x}, \texttt{base\_y}, \texttt{base\_yaw}), not a free-floating joint---same \texttt{xyt} protocol as Stretch with different MuJoCo mechanics.
\end{itemize}

\subsection{Depth policy (critical for experiments)}

MuJoCo already publishes \textbf{rendered sensor depth} on ZMQ.
For sim Dynagraph/DynaMem runs we recommend \texttt{depth\_source: sensor} via the default \texttt{dynav\_config.yaml}.

\textbf{Do not enable Depth Anything 3 in sim} unless deliberately reproducing a DA3-only deployment stack---mixing inferred depth with a sensor-depth pipeline complicates debugging and can desynchronize Rerun point clouds from the voxel map.

On \textbf{real hardware}, the bridge omits \texttt{depth}; clients set \texttt{allow\_missing\_depth} and use \texttt{depth\_source: auto} or \texttt{da3} (see Appendix~\ref{app:da3}).
The maintained preset is \texttt{dynav\_innate\_mars.yaml}; pass \texttt{--dynav-config dynav\_innate\_mars.yaml} for hardware Mars or when matching a DA3 depth path.

\textbf{Sanity modes:}
\begin{itemize}
  \item \texttt{--perfect-depth} / \texttt{EMET\_DYNAMEM\_PERFECT\_DEPTH} forces sensor depth in sim and skips DA3---useful for checking navigation frame, intrinsics, and base pose.
  \item \texttt{emet debug-da3-depth --robot innate\_mars} logs DA3 depth and a strided point cloud to Rerun via the same \texttt{resolve\_depth\_map} path as DynaMem; add \texttt{--depth-source sensor} to compare against sim ground truth.
\end{itemize}

\subsection{Cameras and ZMQ keys}

Primary mapping uses \texttt{head\_left} RGB with optional stereo mate \texttt{rgb\_right} plus matched intrinsics and poses.
Wrist \texttt{camera\_arm} is sim-polished for frustum clearance; the hardware bridge may not publish an arm stream.
Session metadata (\texttt{emet\_session}) reports \texttt{runtime\_kind=innate\_mars\_ros2\_bridge} on hardware.

\subsection{Dynagraph usage}

\begin{verbatim}
# Terminal 1 (sim)
emet serve mujoco --robot innate_mars --headless

# Terminal 2 -- sensor depth in sim (default dynav)
emet run dynagraph --robot innate_mars --robot-ip 127.0.0.1

# Hardware or DA3-matched stack
emet run dynagraph --robot innate_mars --dynav-config dynav_innate_mars.yaml
\end{verbatim}

\subsection{Known DA3 artifacts}

Stereo depth networks can assign finite depth to sky or bright ceiling strips, producing vertical ``wall'' sheets in the voxel map.
\texttt{dynav\_innate\_mars.yaml} sets \texttt{da3\_clip\_max\_m} and \texttt{da3\_ignore\_sky\_fraction\_top} to mitigate this; tune when room lighting or camera tilt leaves phantom obstacles.
